\documentclass{ximera}

\title{Oefeningen op neurale netwerken}
\author{Mila Vervoort}
\license{CC: 0}         % replace with an appropriate license, or set it in xmPreamble

\begin{document}
\begin{abstract}
    Oefeningen op neurale netwerken 
\end{abstract}
\maketitle
\label{xim:Oefeningen op neurale netwerken}

\section*{Oefening 1: Eenvoudig neuraal netwerk}

In de theorie zagen we dat Machine Learning kan worden gebruikt om een verband uit gegevens te leren. In deze oefening bouwen en trainen we
zelf een eenvoudig neuraal netwerk met \texttt{PyTorch}.

\textbf{Gegeven}

We beschikken over gegevens van de vorm $x \longrightarrow y$ en willen een neuraal netwerk trainen dat op basis van $x$ een voorspelling $\hat y$ kan maken.

Voor deze oefening gebruiken we gegevens die afkomstig zijn van de functie $y=x^2+1.$

We doen hierbij alsof deze gegevens afkomstig zijn uit een experiment.
Het neurale netwerk krijgt de onderliggende vergelijking dus niet te zien.
Het krijgt enkel de gegevens en moet daaruit zelf het verband leren.

We gebruiken bijvoorbeeld de volgende waarden van $x$:
\[
-10,-8,-6,-5,-4,-2,-1,0,1,2,4,5,6,8,10.
\]

Voor iedere waarde van $x$ berekenen we de overeenkomstige waarde van $y$.

De berekeningen voer je uit in de bijhorende Google Colab-notebook. Open de notebook via onderstaande link en gebruik deze tijdens het
oplossen van de oefening.

\begin{center}
\href{https://colab.research.google.com/drive/17OMMkgFNeBZeORY04bA9OUqVQtMOB594?usp=sharing}{\Large Open de Google Colab-notebook}
\end{center}

\textbf{Gevraagd}

Gebruik een neuraal netwerk om het verband tussen $x$ en $y$ te leren.
Doorloop hierbij de volgende stappen:

\[
\boxed{
\begin{array}{c}
\text{1. Data voorbereiden}\\
\downarrow\\
\text{2. Neuraal netwerk bouwen}\\
\downarrow\\
\text{3. Netwerk trainen}\\
\downarrow\\
\text{4. Netwerk testen}\\
\downarrow\\
\text{5. Nieuwe waarde voorspellen}
\end{array}
}
\]


\begin{question} \textbf{Stap 1: Data voorbereiden}

Een neuraal netwerk kan niet rechtstreeks met een Python-lijst van getallen werken. De gegevens moeten eerst worden omgezet naar
\textit{PyTorch-tensors}.

Voer de eerste stap in de Google Colab-notebook uit. De trainingsgegevens worden automatisch aangemaakt en vervolgens omgezet naar tensors.

Welke uitspraak is correct?

\begin{multipleChoice}
    \choice{Het neurale netwerk krijgt de vergelijking $y=x^2+1$ als input.}
    \choice{Het neurale netwerk krijgt alleen de waarden van $x$ en moet zelf de bijbehorende waarden van $y$ raden.}
    \choice[correct]{Het neurale netwerk krijgt de trainingsvoorbeelden $(x,y)$ en leert daaruit een verband.}
    \choice{Het neurale netwerk krijgt alleen de exacte waarde van $y$ voor $x=0$.}
\end{multipleChoice}

\begin{feedback} \text{ }

Het neurale netwerk krijgt de trainingsgegevens $(x,y)$ als voorbeeld.
De vergelijking
\[
y=x^2+1
\]
wordt niet aan het netwerk meegegeven.

Het netwerk probeert zelf een functie te leren die de gegeven gegevens zo goed mogelijk beschrijft.
\end{feedback}

\end{question}


\begin{question} \textbf{Stap 2: Bouw het neurale netwerk}

We willen een model dat één waarde $x$ als input krijgt en één waarde $\hat y$ als output geeft:

\[
x
\longrightarrow
f_\theta
\longrightarrow
\hat y.
\]

In deze oefening gebruiken we een netwerk met de structuur

\[
\boxed{1\rightarrow8\rightarrow8\rightarrow1}.
\]

Dit betekent:

\begin{itemize}
    \item één inputneuron;
    \item twee hidden layers met telkens acht neuronen;
    \item één outputneuron.
\end{itemize}

De hidden layers gebruiken de activatiefunctie \texttt{Tanh}.

De structuur van het netwerk is reeds voorgeprogrammeerd in de Google Colab-notebook. Bekijk de code en probeer de verschillende lagen te herkennen. Voer daarna het codeblok uit om het netwerk aan te maken.

We kunnen het netwerk nu al een voorspelling laten maken: $\hat y=f_\theta(x).$

Vul in de Google Colab-notebook de ontbrekende code aan om voor de trainingsgegevens een eerste voorspelling te maken.

Zijn de voorspellingen van het netwerk vóór de training al goed?

\begin{multipleChoice}
    \choice{Ja, want het netwerk kent de vergelijking $y=x^2+1$.}
    \choice[correct]{Nee, de voorspellingen zijn nog niet afgestemd op de gegeven data.}
    \choice{Ja, want PyTorch berekent automatisch de exacte vergelijking.}
    \choice{Nee, omdat een neuraal netwerk geen verbanden kan leren.}
\end{multipleChoice}

\begin{feedback} \text{ }

Het netwerk kent de onderliggende relatie $y=x^2+1$ nog niet. Wanneer het netwerk wordt aangemaakt, krijgen de parameters zoals de gewichten en biases beginwaarden. Deze waarden zijn nog niet afgestemd op onze trainingsgegevens.

Het netwerk kan dus al wel een voorspelling maken, $\hat y= f_\theta(x),$ maar deze voorspelling zal in het begin nog niet goed overeenkomen met de werkelijke waarden.

Tijdens het trainen worden de parameters $\theta$ aangepast zodat de
voorspellingen steeds beter worden.

\end{feedback}

\end{question}

\begin{remark} \text{ }
    
De keuze voor twee hidden layers met telkens acht neuronen is hier vooral een praktische keuze. Het verband $y=x^2+1$ is relatief eenvoudig, waardoor een klein netwerk voldoende is om het verband te leren. Er bestaat geen algemene regel die zegt dat voor dit probleem precies acht neuronen of twee
hidden layers nodig zijn.

Bij andere problemen kan de optimale structuur van een neuraal netwerk anders zijn. Het aantal lagen en neuronen moet dan vaak worden getest en geoptimaliseerd. Een te klein netwerk kan onvoldoende complexe verbanden leren, terwijl een onnodig groot netwerk meer parameters
heeft en moeilijker te trainen kan zijn.

\end{remark}

\begin{question} \textbf{Stap 3: Het netwerk trainen}

Zoals we in de vorige stap hebben gezien, zijn de voorspellingen van het ongetrainde netwerk nog geen goede benadering van onze gegevens. We willen nu de parameters van het netwerk aanpassen zodat de voorspellingen steeds beter worden.

Hiervoor hebben we twee dingen nodig:

\begin{enumerate}
    \item een manier om te bepalen hoe groot de fout van het netwerk
    is;
    \item een methode om de parameters van het netwerk aan te passen.
\end{enumerate}

De eerste stap gebeurt met behulp van een \textbf{lossfunctie}. De tweede stap gebeurt met behulp van een \textbf{optimizer}.

Voor dit probleem gebruiken we de \textit{Mean Squared Error (MSE)} als lossfunctie en het \textit{Adam-algoritme} als optimizer.

Voer de code voor het definiëren van de loss-functie en optimizer uit in de Google Colab-notebook. 

We hebben nu alles wat nodig is om het netwerk te trainen. Tijdens één volledige trainingsstap, ook wel een \textbf{epoch} genoemd, gebeurt het volgende:

\[
\boxed{
\begin{array}{c}
\text{oude gradiënten verwijderen}\\
\downarrow\\
\text{nieuwe voorspelling maken}\\
\downarrow\\
\text{loss berekenen}\\
\downarrow\\
\text{backpropagation}\\
\downarrow\\
\text{parameters aanpassen}
\end{array}}
\]

Dit proces wordt vervolgens vele malen herhaald.

In de Google Colab-notebook is de trainingsloop al voorgeprogrammeerd. Bekijk de verschillende stappen in de code en voer de volledige trainingsloop uit.

De training wordt hier $10000$ keer herhaald. Tijdens de training wordt de loss van iedere epoch opgeslagen in \texttt{loss\_history}. We kunnen deze waarden gebruiken om te bekijken hoe de loss verandert tijdens het trainen.

Schrijf in de Google Colab-notebook een code die de \texttt{loss\_history} uitzet tegen het aantal epochs. Bekijk vervolgens de bekomen grafiek. 

Wat kan je uit de grafiek besluiten over het leerproces?

\begin{multipleChoice}
    \choice{De loss stijgt in het algemeen, \\
    wat betekent dat de gemiddelde kwadratische fout tijdens de training groter wordt en het netwerk de trainingsdata steeds beter benadert.}
    \choice{De loss daalt, \\
    dus de voorspellingen van het netwerk worden bij
    elke epoch voor alle afzonderlijke trainingspunten beter.}
    \choice[correct]{De loss daalt in het algemeen, \\
    wat erop wijst dat de gemiddelde kwadratische fout op de trainingsdata kleiner wordt tijdens
    het trainen.}
    \choice{De loss daalt, \\
    dus het netwerk heeft na de laatste epoch de
    exacte relatie $y=x^2+1$ geleerd en zal voor elke nieuwe input de ouput exact voorspellen.}
\end{multipleChoice}

\begin{feedback} \text{ }

Een mogelijke code is:

\begin{verbatim}
plt.plot(loss_history)
plt.xlabel("Epoch")
plt.ylabel("Loss")
plt.title("Loss tijdens de training")
plt.show()
\end{verbatim}

De grafiek toont hoe de loss verandert terwijl het netwerk wordt getraind.

Tijdens het trainen daalt de loss in het algemeen. Dit betekent dat de voorspellingen van het netwerk steeds beter overeenkomen met de trainingsdata.
\end{feedback}
\end{question}


\begin{question} \textbf{Stap 4: Het netwerk testen}

Het netwerk is nu getraind. We willen nagaan of het geleerde model het verband tussen $x$ en $y$ inderdaad goed benadert. Daarom laten we het getrainde netwerk nu voorspellingen maken voor een hele reeks waarden van $x$.

We gebruiken hiervoor 300 gelijkmatig verdeelde waarden tussen $-10$ en $10$.  

We willen nu voor elke $x$-waarde een voorspelling van het getrainde netwerk berekenen. Vul in Google Colab de ontbrekende code aan, zodat het netwerk voor alle waarden in \texttt{x\_plot} een voorspelling \texttt{y\_plot} genereert. \\
Bekijk vervolgens de plot waarin de oorspronkelijke trainingsdata wordt weergegeven samen met de voorspellingen van het neurale netwerk en vergelijk.

Welke conclusie past het best bij een grafiek waarin de rode curve de trainingspunten goed volgt?

\begin{multipleChoice}

\choice{Het netwerk heeft de trainingsdata opgeslagen, maar heeft geen
verband tussen $x$ en $y$ geleerd.}

\choice[correct]{Het netwerk heeft een functie geleerd die de relatie tussen
$x$ en $y$ in het bereik van de trainingsdata goed benadert.}

\choice{Omdat de rode curve de trainingsdata goed volgt, is bewezen dat het
netwerk voor elke mogelijke waarde van $x$ exact $x^2+1$ zal voorspellen.}

\choice{Het netwerk voorspelt alleen de trainingspunten en kan geen waarden voor $x$ voorspellen die niet in de trainingsdata voorkomen.}

\end{multipleChoice}

\begin{feedback} \text{ }

Wanneer de rode curve de trainingspunten goed volgt, betekent dit dat het getrainde netwerk een functie heeft geleerd die de relatie tussen $x$ en $y$
in dit gebied goed benadert.

Het netwerk probeert dus niet letterlijk de tabel met waarden te onthouden, maar leert een functie

\[
\hat y=f_\theta(x)
\]

waarmee het ook voorspellingen kan maken voor waarden van $x$ die niet letterlijk in de trainingsdata stonden.

We moeten hierbij wel voorzichtig zijn met onze conclusie. Een goede overeenkomst met de trainingsdata betekent niet automatisch dat het netwerk voor iedere mogelijke nieuwe waarde van $x$ een perfecte voorspelling zal
geven. \\
Vooral buiten het bereik van de trainingsdata kan een neuraal netwerk
slecht extrapoleren.
\end{feedback}

\end{question}

\begin{question} \textbf{Stap 5: Een nieuwe waarde voorspellen}

Een belangrijk doel van Machine Learning is dat een getraind model ook voorspellingen kan maken voor nieuwe input.

Tijdens de training heeft het netwerk geleerd om een verband te benaderen tussen $x$ en $y$. We kunnen het getrainde model nu gebruiken om voor een nieuwe waarde van $x$ een voorspelling te maken.

We kiezen bijvoorbeeld een waarde die ook binnen het bereik van onze trainingsdata $[-10, 10]$ ligt: $x_{\text{new}}=7.$ \\
(Je zou ook een andere waarde kunnen controleren die binnen het bereik van onze trainingsdata ligt.)

Vul vervolgens de code aan zodat het getrainde netwerk een voorspelling maakt voor deze nieuwe waarde. \\
Vergelijk de voorspelling met de exacte waarde van $y=x^2+1.$ Welke conclusie is dan het meest correct?

\begin{multipleChoice}

\choice{Het netwerk heeft niet geleerd, want een neuraal netwerk moet altijd
de exacte waarde voorspellen.}

\choice{Het netwerk heeft de trainingsdata fout opgeslagen, want de
voorspelling moet exact gelijk zijn aan $50$.}

\choice[correct]{Het netwerk heeft de relatie goed benaderd, maar de
voorspelling wijkt nog enigszins af van de exacte waarde.}

\choice{De voorspelling is correct, want $49.2$ en $50$ zijn dezelfde waarde
binnen een neuraal netwerk.}

\end{multipleChoice}

\begin{feedback} \text{ }

De exacte waarde is $y=7^2+1=50,$
terwijl het netwerk bijvoorbeeld $49.2$ voorspelt. De voorspelling is dus niet exact, maar ligt wel dicht bij de werkelijke waarde.

Een neuraal netwerk leert de relatie uit de beschikbare trainingsdata. Het doel van de training is daarom niet noodzakelijk om voor elk punt exact de juiste waarde te verkrijgen, maar om een functie te vinden die de gegevens zo goed mogelijk benadert.

De nauwkeurigheid van het netwerk kan verder worden verbeterd door het model te optimaliseren. Dit kan bijvoorbeeld door meer of betere trainingsdata te gebruiken, de structuur van het netwerk aan te passen (zoals het aantal lagen of neuronen), of de trainingsinstellingen, zoals de learning rate en het aantal trainingsrondes (epochs), te wijzigen.

\end{feedback}

\end{question}



\end{document}